% =====================================================================
%  Financial Networks Project: working notes
%  Running log. Add a new dated section at the bottom as work happens.
%  Compiles on Overleaf as-is; no .bib, no images.
% =====================================================================
\documentclass[11pt,a4paper]{article}

\usepackage[T1]{fontenc}
\usepackage[utf8]{inputenc}
\usepackage{lmodern}
\usepackage[margin=2.3cm]{geometry}
\usepackage{amsmath,amssymb}
\usepackage{booktabs}
\usepackage{enumitem}
\usepackage{xcolor}
\usepackage{microtype}

\definecolor{BlueLink}{HTML}{1F4E79}
\definecolor{BadRed}{HTML}{B3261E}

\usepackage[colorlinks=true,linkcolor=BlueLink,urlcolor=BlueLink]{hyperref}
\emergencystretch=3em

% Comments: use Overleaf's own review mode (Review -> Add comment).

\newcommand{\Q}{Q}
\newcommand{\Lap}{\mathbf{L}}
\newcommand{\Amat}{\mathbf{A}}
\newcommand{\Bmat}{\mathbf{B}}
\newcommand{\Cmat}{\mathbf{C}}
\newcommand{\logday}[1]{\par\vspace{3mm}\noindent{\large\textbf{#1}}\par\vspace{1mm}\hrule\vspace{2mm}\noindent}

\setlist[itemize]{topsep=2pt,itemsep=1pt,parsep=0pt,leftmargin=1.3em}
\setlist[enumerate]{topsep=2pt,itemsep=1pt,parsep=0pt,leftmargin=1.6em}
\pagestyle{plain}
\setlength{\parskip}{2pt}

\begin{document}

\begin{center}
{\large\textbf{Financial Networks Project: working notes}}\\[1mm]
{\small Jash Dalal, last updated 8 September 2026. Newest entries at the bottom.}
\end{center}

\logday{5 July: Paper 1, learning a Laplacian from returns}

Builds a network of stocks straight from price movements instead of building a
correlation table and cutting it up. Watches one number, $\lambda_2$, which says
how tightly the network holds together; small means close to breaking apart, and
the rule is to stay invested below a set level and sit in cash above it.

The method, the month-to-month movement of $\lambda_2$, and the COVID dodge all
reproduce. The money does not. Here $\lambda_2$ stays high through April 2020
where the paper's drops, so the strategy sits in cash through the recovery and
ends at $-0.011$ where simply holding ends at $+0.078$. No threshold was moved
to fix it. Takeaway: $\lambda_2$ is the incumbent indicator on the Laplacian
side, and it is fragile.

\logday{22 to 28 July: Paper 2, market regimes with TICC}

Cuts time into stretches, each with its own pattern of which series depend on
which, and the same pattern can come back later. Headline accuracy reproduces.
The side experiments that switch parts of the method off come out weaker than
reported, so this rebuild makes the method look \emph{more} necessary than the
paper's own numbers do. Peripheral to the current question; it is the only one
of the three that models a state stretched over time, so keep it as the
comparator if a network indicator ever needs checking against a plain
regime model.

\logday{9 to 10 August: Paper 3, modularity of market networks}

Each day, keep the strongest 10\% of pairwise correlations as connections, run
Louvain, record one number $\Q$ for how cleanly the stocks split into groups.
Rebuilt at full size: 360 stocks, 25 years, 6{,}315 daily networks. It holds up
on its own terms. The trouble is what the number means.

\begin{itemize}
\item \textbf{$\Q$ does not show the market has groups.} Pure random data
through the same pipeline gives $0.2246$; the real market gives $0.2215$. Noise
comes out marginally \emph{more} grouped than the stock market. Cutting up a
correlation table always yields a network full of triangles, and those always
look grouped.
\item \textbf{Louvain disagrees with itself.} Run-to-run spread on the same day
is $2.09\times$ a normal day-to-day move in $\Q$, on 78\% of days.
\item \textbf{The discarded number works better.} $\tau(t)$, the cut-off needed
each day to keep 10\% of pairs, is computed, used, then thrown away. It scores
$0.807$ against $\Q$'s $0.538$.
\item \textbf{Compression fails where the story is.} The group summary is
cheaper than the network on 69\% of ordinary days but only 34\% of crisis days,
because in a crisis the network shatters and the summary outgrows the thing it
was summarising.
\end{itemize}

Re-audit the same day: the paper's four non-reproducing numbers are real and
survive nine implementation variants, but the cause is that these networks fall
into disconnected pieces on 63\% of days (99.3\% after 2002). Eight defects were
found in this rebuild's own criticisms and fixed.

\logday{12 August: the two matrices behind the two indicators}

\textbf{The prompt.} Paper 1 reads $\lambda_2$ off the graph Laplacian. Paper 3
reads $\Q$ off a grouping of stocks. Modularity has a matrix underneath it that
nobody in this literature reports. Is there a study in comparing the two?

\textbf{Laplacian} $\Lap = \mathbf{D} - \Amat$: how well connected the network
is and where it would break. Number of zero eigenvalues is how many separate
pieces it has fallen into; $\lambda_2$ is zero when it is already in pieces and
large when solidly joined.

\textbf{Modularity matrix} $\Bmat = \Amat - \mathbf{d}\mathbf{d}^{\!\top}/2m$:
what was seen minus what was expected. The baseline $d_i d_j/2m$ answers ``how
many connections would these two have by accident, given how popular each is''.
Two typical stocks (36 connections each) give $0.10$, exactly the network's
density. Two heavily connected stocks (100 each) give $0.77$, so a link between
them is not news. Every entry is a surprise score:

\begin{center}\small
\begin{tabular}{@{}llll@{}}
\toprule
seen & expected & entry & reading \\
\midrule
linked & $0.10$ & $+0.90$ & together far more than chance \\
linked & $0.77$ & $+0.23$ & two popular stocks linked; barely news \\
not linked & $0.10$ & $-0.10$ & ordinary; most pairs look like this \\
not linked & $0.77$ & $-0.77$ & two popular stocks \emph{avoiding} each other \\
\bottomrule
\end{tabular}
\end{center}

That last row is the point: $\Bmat$ has a meaningful negative end, while the
Laplacian, built only from connections that exist, cannot say ``should be there
and is missing''. And $\Q$ is just a sum of selected entries of $\Bmat$, so the
chain runs from a $360\times360$ matrix, through one grouping picked by an
unstable search, down to a single number. Paper 3 reports the last box.

Reading $\Bmat$ directly instead: $\mu_1$ measures the strongest split (same job
as $\Q$, but deterministic), $\#\{\mu_i>0\}$ caps the number of separable groups
so it is a group count with no algorithm involved, and $\mu_N$ picks out
\emph{anti}-groups, sets of stocks avoiding each other. None need Louvain, and
none appear anywhere in finance.

\logday{12 August: spectral probe over the full 25 years}

Full spectrum of $\Lap$, $\Bmat$ and the raw correlation table $\Cmat$ for all
6{,}315 days, every statistic through the same detector: a 250-day trailing
$z$-score using only prior days, scored against crisis dates fixed before any of
this work. Score is AUC, the chance a random crisis day ranks above a random
calm day, so $0.50$ is a coin flip. Six minutes on one core.

\begin{center}\small
\begin{tabular}{@{}llr@{}}
\toprule
\textbf{Signal} & \textbf{What it measures} & \textbf{AUC} \\
\midrule
spectral entropy of $\Cmat$ & how spread out the correlation structure is & \textbf{0.822} \\
$\lambda_1(\Cmat)/N$ & share of movement in the biggest common factor & 0.812 \\
$\tau(t)$ & cut-off needed that day to keep 10\% of pairs & 0.807 \\
$\mu_N(\Bmat)$ & \textbf{strongest anti-grouping} & \textbf{0.760} \\
isolated-node count & stocks left with no connections & 0.756 \\
$\sum_{\mu_i>0}\mu_i$ & \textbf{total group mass} & 0.750 \\
$\lambda_N(\Lap)$ & largest Laplacian eigenvalue & 0.731 \\
$\#\{\mu_i>0\}$ & \textbf{group count, no algorithm} & 0.722 \\
$\lambda_2(\Lap)$, largest piece & Fiedler value, breakage removed & 0.652 \\
$\mu_1-\mu_2$ & modularity spectral gap & 0.636 \\
\# modes above the noise band & how many groups are even resolvable & 0.619 \\
$\lambda_2(\Lap)$, whole network & Fiedler value, \textbf{Paper 1} & 0.617 \\
$\mu_1(\Bmat)$ & leading modularity eigenvalue & 0.607 \\
$\Q(t)$ & Louvain modularity, \textbf{Paper 3} & \textbf{0.538} \\
$\max_k(\lambda_{k+1}-\lambda_k)$ & largest Laplacian gap & 0.526 \\
\bottomrule
\end{tabular}
\end{center}

\textbf{1. The matrix beats the single number.} $\Q$ at $0.538$ means it ranks a
crash day above a calm day 53.8\% of the time; a coin manages 50\%. Three
separate readings of the matrix $\Q$ comes from reach $0.72$ to $0.76$ on the
same networks under the same test, so the information was in $\Bmat$ all along
and the Louvain step is where it is lost. The best modularity signal is the
\emph{negative} end, which thresholding structurally cannot see. And $\mu_1$
tracks $\Q$ at $r = 0.761$ while being fully deterministic, so it is a drop-in
with none of the wobble. Paper 1's $0.617$ is not a refutation of Paper 1, which
learns a regularised Laplacian on 7 stocks, a different object entirely.

\textbf{2. The identity breaks, and that is the good news.} On a network where
every stock has the same number of connections $d$, it is exactly true that
$\lambda_2(\Lap) + \mu_1(\Bmat) = d$: the Fiedler value and the leading
modularity eigenvalue are the same number read backwards. If that held here,
comparing Paper 1's indicator with Paper 3's would be comparing a thing to
itself. Paper 3's construction pins the average degree at exactly $35.90$ every
day, so it nearly should hold. Measured over all days the sum falls short by
$\mathbf{9.49 \pm 4.63}$, i.e.\ 26\% of the average degree, and the $\pm4.63$
says it is not a fixed offset that could be subtracted, it moves by a third of
its own size day to day. Restricting to the 2{,}342 one-piece days gives
$9.21 \pm 3.60$; the largest piece only gives $11.11 \pm 5.31$. Cleaner still:
two quantities summing to a constant would be perfectly anti-correlated at $-1$,
unrelated ones give $0$, and the measured value is $\mathbf{-0.187}$. Dual in
theory, unrelated in practice. So the question stops being ``which indicator
wins'' and becomes \emph{what is that $9.49$ made of}: uneven degrees, breakage,
or the market-wide component.

\textbf{3. There may be nothing to detect.} 64{,}620 correlations are estimated
from 30 days on 360 stocks, six times more output than input, and at most 29 of
the 360 eigenvalues can be non-zero by arithmetic alone. Random-matrix theory
puts the noise band's upper edge at $19.93$, so anything below that is
indistinguishable from what pure noise produces. Counting eigenvalues that
escape it, excluding the market-wide component: \textbf{mean $0.70$, max $3$,
often zero}. Less than one resolvable group per day, while Louvain confidently
reports a median of nine and the paper narrates their 25-year dynamics. This
indicts the whole literature, not just Paper 3; a 30-day window on a few hundred
stocks is standard.

\textbf{4. The strong signals are one signal, and it does not generalise.}
Spectral entropy, the absorption ratio and $\tau(t)$ move together at $|r| >
0.99$, so they are the same series with different labels, all measuring the
market-wide component. That is not four discoveries but one, made in 1999.
Training before 2002 and testing after (27 crisis days in the test period):

\begin{center}\small
\begin{tabular}{@{}lrr@{}}
\toprule
& training years & unseen years \\
\midrule
market-wide component alone & $0.812$ & $\mathbf{0.730}$ \\
\quad $+\ \mu_1(\Bmat)$ & $0.822$ & $0.629$ \\
\quad $+\ \mu_N(\Bmat)$ & $0.810$ & $0.719$ \\
\quad $+\ \Q(t)$ & $0.831$ & $0.699$ \\
\quad $+\ \lambda_2(\Lap)$, largest piece & $0.810$ & $0.700$ \\
\bottomrule
\end{tabular}
\end{center}

Every addition helps on the training years and hurts on the unseen ones, the
textbook overfitting fingerprint. 27 events is too few to settle it either way,
but ``too weak to settle'' is not a defence. One bright spot: $\mu_1$ tracks the
market-wide component at only $0.025$ while tracking $\Q$ at $0.761$, so it
really is a new axis. $\mu_N$ tracks it at $0.68$, so much of its $0.760$ is
borrowed, which is exactly the D3 gate below.

Code and outputs in \texttt{reports/probe/}. It imports the rebuild's own crisis
dates, network construction and scoring unchanged; as a check, $\Q$ and $\tau$
come out at $0.5380684370$ and $0.8071316467$, matching the stored values to ten
places.

\logday{12 August: proposed directions and open questions}

\textbf{D1, what breaks the identity.} Decompose the $9.49$ into uneven degrees,
breakage, and the market-wide component, and see which carries the stress
information. Same spectrum over a family of matrices ($\Amat$, $\Lap$,
normalised $\Lap$, $\Bmat$, normalised $\Bmat$, and the threshold-free
correlation modularity of MacMahon \& Garlaschelli) under one shared test. Not a
bake-off but a structural result with a bake-off attached. Risk: the gap may be
dominated by breakage, an artifact of thresholding. The largest-piece variant
already survives, but that needs establishing across constructions.

\textbf{D2, is there anything to detect.} Map, over window length, universe size
and density, where the market sits relative to the noise band and the
Nadakuditi and Newman detectability threshold, so anyone in this literature can
check whether their settings are inside the region where grouping means
anything. Ship a deterministic replacement for the group count. Risk: if
essentially no configuration in the literature is inside that region, the
contribution is purely negative. Still publishable, harder to place.

\textbf{D3, the negative end (gated).} Paper 3 discards $4{,}246{,}798$
pair-days of strong \emph{anti}-correlation, recorded identically to ``no
relationship'', and that series collapses to near zero from about 2003 and never
recovers, right beside the paper's own unexplained 2002 break which it lists as
future work. $\mu_N$ at $0.760$ is the best new number here. Gate before
committing: does $\mu_N$ survive removing the market-wide component? Two-day
check.

\textbf{Recommendation: D1 and D2 as one project, D3 second if the gate passes.}
D1 alone risks being a bake-off, D2 alone risks being purely negative. Together
they are one claim: here is the right matrix to read market group structure
from, here is where reading it means anything at all, and here is what the
field's two incumbent indicators actually measure.

\textbf{Open questions.}
\begin{enumerate}[label=\textbf{Q\arabic*.}]
\item Is the identity the right spine, or would you prefer the plainer empirical
comparison?
\item Keep Paper 3's thresholded networks (comparable to the rebuild) or move to
threshold-free correlation modularity (more defensible, breaks comparability)?
Inclination is both, and report the difference, since that difference is the
``beyond thresholding'' argument made concrete.
\item 27 unseen crisis days is far too few. Extend to 2026, add exchanges, or go
index-level?
\item Venue: a physics venue wants the spectral and detectability result, a
finance venue wants the indicator to beat the absorption ratio on unseen data,
which on current evidence it does not.
\item Carry TICC forward as a comparator, or park it?
\end{enumerate}

\logday{19 August: the exploration written up properly}

The 12 August entries above are too compressed to follow, so the whole
exploration is now written out in \texttt{reports/2026-08-19\_discrepancy\_explained.tex} and
shared for comment. It builds every result up from the papers' own definitions:
the three replications side by side, each observed discrepancy beside what would
explain it, the one scoring apparatus everything is judged on, and the three
directions in a single symmetric frame. 17 pages, eight diagrams,
self-contained, Overleaf-ready; every number in it is re-derived from the stored
artifacts by \texttt{reports/probe/check\_numbers.py}.

Three figures quoted above are stated more carefully there. The $0.68$ and
$0.025$ market-mode correlations hold on the $z$-scored detector series, not in
levels (where they are $-0.819$ and $+0.294$); the AUC-$0.756$ row is the count
of \emph{connected components}, not of isolated nodes; and the noise-control
$\Q$ values come from two different runs and are not interchangeable.

\logday{31 August: the meeting, and the question set}

Asked to take the modularity matrix seriously as a source of crisis indicators:
find a quantity derived from $\Bmat$ that carries the same information as
$\lambda_2(\Lap)$, tested on Paper 1's own FAAMUNG data, so the two can be
declared equivalent and a theory framework built on top. Not necessarily a
second-smallest eigenvalue, anything at all. Scope agreed: an empirical menu, Paper 1's
dataset only, no Paper 3 scoring.

Also raised at the meeting: the $\lambda_2(\Lap) + \mu_1(\Bmat) = d$ spine
proposed as D1/D2 on 12 August holds only when \emph{every} degree equals $d$.
Substituting a mean degree is not a weakened hypothesis, it is no hypothesis.
\textbf{The 12 August framing of D2 is retracted}; the quantitative confirmation is
in the next entry.

\logday{31 August to 7 September: a modularity equivalent of the Fiedler value}

Paper 1 throws the learned graphs away after reading $\lambda_2$ off them, so
Algorithm 2 was re-run with its exact constants to recover all 243 Laplacians,
$\Amat_t = -\operatorname{offdiag}(\Lap_t)$ recovered from each, and 22
modularity-derived quantities scored against 6 Laplacian references. Gates first:
recomputed $\lambda_2$ matches the cached series to $1.2\times10^{-10}$, and the
backtest reproduces the published $S1 = +0.078$ against $S2 = -0.011$.

\textbf{The question dissolves under normalisation.} Writing
$\mathbf{u} = \mathbf{D}^{1/2}\mathbf{1}/\sqrt{2m}$, the normalised modularity
matrix is $\mathbf{D}^{-1/2}\Bmat\mathbf{D}^{-1/2} =
\mathbf{I} - \Lap_{\mathrm{sym}} - \mathbf{u}\mathbf{u}^{\!\top}$, and $\mathbf{u}$
is $\Lap_{\mathrm{sym}}$'s own null vector. The two normalised matrices are the
same operator: eigenvalues paired one-to-one by $\mu = 1-\lambda$ (full spectrum,
all 243 windows, $1.8\times10^{-15}$) and eigenvectors shared to $1.000000$. So no
normalised-$\Bmat$ indicator can be new; it is a Laplacian indicator in different
notation. Unnormalised, the two do \emph{not} share a basis (overlap $0.871$; on
some windows the Fiedler vector and $\Bmat$'s leading mode are nearly orthogonal,
$0.039$), so that is where anything new must live.

\textbf{The best-scoring candidate was the spurious one.} $\bar{d}-\mu_1(\Bmat)$
reproduced 91\% of $\lambda_2$'s daily invest/cash calls, and it is exactly the
$\bar d$ substitution flagged at the meeting. Whole-spectrum error $1.25$, i.e.\
$5.3\times$ $\lambda_2$'s own sd; residual sd $0.208$, \textbf{89\% of $\lambda_2$'s
sd}; only 40\% of the proxy's variance is $\lambda_2$; and the residual is
\emph{smaller} in-crash ($0.308$ vs $0.512$), so the error term was doing some of
the detecting. Degree CV here is $0.305$, nowhere near regular. The $\bar d$ term
also does no work: controlling for it, $\mu_1(\Bmat)$ alone gives $-0.851$ against
the proxy's $+0.851$.

\textbf{Correlation was the wrong screen.} It was the working comparison metric
between $\Bmat$ and $\Lap$, and it fails both ways. It promoted the spurious
candidate above; and spectral entropy of the normalised $\Bmat$ correlates only
$-0.264$ with $\lambda_2$, which looks independent, yet rebuilds from the
$\Lap_{\mathrm{sym}}$ spectrum alone to $8.0\times10^{-15}$. Novelty has to mean
\emph{not reconstructible from} $\operatorname{spec}(\Lap)$, not \emph{decorrelated
from} $\lambda_2$.

\textbf{$\lambda_2$ is itself inside the noise band on Paper 1's own setup.} Seven
stocks, 243 windows, one event, 24 crash days. Raw-level COVID AUC is
$\mathbf{0.472}$, below chance; only a strictly-prior $z$-score lifts it to
$0.759$. It clears a bootstrap over days (CI $[0.634,0.875]$) but fails the harder
null: sliding a 24-day window anywhere in the series puts COVID at the
\textbf{76th percentile, $p = 0.243$}. A random gate at the same invest rate beats
buy-and-hold 39.2\% of the time, so $\sim$11 of 28 candidates were expected to beat
$S1$ by luck; 5 did, none with $p < 0.115$. No multiplicity correction exists
anywhere in this study.

\textbf{What survives, and it reframes the project.} $\lambda_2$ is read off the
\emph{unnormalised} $\Lap$, so it is a blend: $d_{\min}\lambda_2(\Lap_{\mathrm{sym}})
\le \lambda_2(\Lap) \le d_{\max}\lambda_2(\Lap_{\mathrm{sym}})$ holds on 243/243
windows. Splitting the detection: structure alone $0.700$, scale alone $0.537$ (a
coin flip), the two combined $0.759$. \textbf{Degree information is a modifier, not
a signal}: worthless alone, but it improves the signal when mixed in, and it is
exactly what normalisation destroys. $\Lap$ combines it additively on the diagonal;
$\Bmat$ combines it through a random-rewiring null. That difference is the only
thing $\Bmat$ offers, and it is the question worth taking forward. It also explains
rather than merely observes why $\mu_N(\Bmat)$ keeps winning: $z$-AUC $0.775$ here
against $\lambda_2$'s $0.759$, redundancy only $0.51$, and independently the best
modularity signal in \texttt{reports/probe/} ($0.760$ vs $0.617$ over 17 crises).
Not significant on one event ($\Delta$AUC $+0.015$, CI $[-0.115,+0.148]$, and
$\mu_N$ fails the same sliding null at $p = 0.345$), so it is a lead, not a result.

Written up in full in \texttt{reports/2026-09-08\_modularity\_equivalence.tex}
(10 pages);
code and the number-checking script in \texttt{modularity\_indicator/}.

\logday{\textcolor{gray}{dd Month: next entry here}}

\vspace{4mm}
{\footnotesize
\textbf{Pointers.} Papers: arXiv:2005.09958 (Cardoso \& Palomar),
arXiv:1706.03161 (TICC), arXiv:1501.05040 (Silva et al.). Key outside
references: Kritzman et al.\ 2011 (absorption ratio, the baseline to beat);
Newman 2006 (modularity matrix); Fasino \& Tudisco 2014/2016/2017 (spectrum,
group counts, anti-groups); Nadakuditi \& Newman 2012 (detectability); MacMahon
\& Garlaschelli 2015 (the correlation-specific null Paper 3 needed and did not
use); Masuda et al., \emph{Physics Reports} 2025 (the ``beyond thresholding''
review); Kang et al.\ 2025 (largest Laplacian gap).
\textcolor{BadRed}{Two Laplacian-energy papers (Eur.\ J.\ Finance 29(9) 2023;
Int.\ Rev.\ Econ.\ Finance 2025) have unverified author lists, do not cite
as-is.}
In this repo: \texttt{reports/archive/EXPLORATION\_EXPLAINED.md} for the ground-up
version of all of the above, \texttt{reports/probe/} for code and outputs,
\texttt{modularity\_indicator/} for the 31 August to 7 September work and
\texttt{reports/2026-09-08\_modularity\_equivalence.tex} for its write-up.
Paper 3 needs \texttt{/opt/anaconda3/bin/python} and \texttt{python-igraph}.
\par}

\end{document}
